\documentclass[10pt,a4paper]{amsart}
\usepackage[margin=19mm]{geometry}
\usepackage{amsmath,amssymb,mathtools}
\usepackage[scr=rsfs,cal=euler]{mathalfa}
\usepackage{xfrac}
\usepackage[unicode,hidelinks,bookmarks=false]{hyperref}
\newcommand{\ie}{i.e.\ }
\newcommand{\cf}{cf.\ }
\newcommand{\resp}{respectively\ }
\newcommand{\Subsection}{\subsection}
\providecommand{\nobreakdash}{\nobreak\hbox{-}}
\setlength{\emergencystretch}{2em}
\multlinegap0pt
\usepackage{amssymb}
\newtheorem{thm}{Theorem}
\DeclareMathOperator{\ac}{ac}
\DeclareMathOperator{\cl}{cl}
\title{Evil Twins in Sums of Wildflowers}
\author{Simon Rubinstein-Salzedo, Stephen Zhou}
\date{Source: arXiv:2603.12225v1 (CC BY 4.0)}
\begin{document}
\maketitle

\noindent\emph{Source and licence.} Evil Twins in Sums of Wildflowers. Version arXiv:2603.12225v1 (CC BY 4.0), distributed by arXiv under the Creative Commons Attribution 4.0 International licence, http://creativecommons.org/licenses/by/4.0/, as recorded in the arXiv licence metadata for this exact version and shown on the abstract page https://arxiv.org/abs/2603.12225v1. Full licence notice: \texttt{licenses/arxiv-2603.12225-CC-BY-4.0.txt}.\par\smallskip

\noindent\textit{Editorial scope.} Counted unit: the article's thm (a-k) (source0.tex lines 297--306). The article's complete original proof of it is delivered with it (source0.tex lines 307--324, one \texttt{proof} environment of 18 lines). No mathematics is written, repaired or re-derived: the statement and the proof are the article's own lines, in the article's own notation, and no retained line is substituted or re-flowed.  No further source-local context is needed: every cross-reference of every retained line resolves inside the two delivered spans.  Nothing was omitted from any retained span, so the package carries no omission record.  No retained line contains a citation, so the wrapper carries no bibliography and no bibliography entry.  No retained line contains a figure environment, an image-inclusion command or drawing code, so no image is delivered and the package declares no asset dependency.  Editorial formatting only, in addition: an AMS \texttt{amsart} preamble that restates the article's own declarations and macros used by the retained lines, namely the environments \texttt{thm} on separate counters, together with the standard packages those lines need.  The retained environments print in this document as Theorem 1 in order of appearance, whereas the article prints the same items with the numbering of its own full document.  The complete native source of the article as distributed in the arXiv e-print for this exact version is bundled unchanged as \texttt{sources/196266/source0.tex}.
\begin{thm}
\label{add to a- (a-k)}
   Let $(\mathscr{A},\mathscr{K})$ be dicotic and evilly normal with the property that $A^\mathcal{L} \cap (\mathscr{A}\setminus\mathscr{K})$ and $A^\mathcal{R} \cap (\mathscr{A}\setminus\mathscr{K})$ are star-closed for any $A \in \mathscr{K}$ and that $A^{\mathcal{L}}\cap \mathscr{K}$ and $A^{\mathcal{R}}\cap \mathscr{K}$ are star-closed for any $A \in (\mathscr{A} \setminus \mathscr{K})$.
    
Let $\mathscr{B}$ be a set of games such that $\mathscr{A} \cup \mathscr{B}$ is hereditarily closed. Assume that for all $B \in \mathscr{B}$, $B^\mathcal{L{}} \cap \cl(\mathscr{K} \cup \mathscr{B})$ and $B^\mathcal{L{}} \cap \cl(\mathscr{K} \cup \mathscr{B})$ are star-closed.

Furthermore, assume that if either player wins by going first in $G+*$ to $G$ for any $G\in \ac( (\mathscr{A}\setminus \mathscr{K} )\cup \mathscr{B})) $ under either play convention, there exists another winning move for that player in $G+*$ in the same play convention.
   
   Then $(\cl(\mathscr{A}\cup \mathscr{B}), \ac(\mathscr{K{ \cup \mathscr{B}}}))$ is evilly normal.
\end{thm}

\begin{proof}
    The elements of $\cl(\mathscr{A} \cup \mathscr{B})$ have form $G= A+B$, where $A \in \mathscr{A}$ and $B \in \mathscr{B}$. We will induct on the birthday of the normal-play canonical form of $G$. This is stronger than induction on the literal form of $G$ since we can assume the theorem is true for any game equal in normal play to an option of $G$.

    Say that $A \in \mathscr{K}$. Then we need to show that $o^+(G) = o^-(G)$. Assume that a player, say Left, wins $G$ going first in normal play. We will show that Left also wins $G$ going first under mis\`ere play.

Say that Left wins $G$ going first under normal play by moving to $G^L$, where $G^L \in \ac(\mathscr{K{ \cup \mathscr{B}}})$. Then by induction, $o^+(G^L) = o^-(G^L)$, so she wins $G$ going first under mis\`ere play with the same move. If Left wins $G$ going first under normal play by moving to $G^L \in \ac( (\mathscr{A}\setminus \mathscr{K} )\cup \mathscr{B})$, then it must be that $G^L = A^L+B$ for some $A^L \in \mathscr{A}\setminus \mathscr{K} $. Since $A^{\mathcal{L}} \cap (\mathscr{A} \setminus \mathscr{K})$ is star-closed, there exists a left option ${A^L}' \in \mathscr{A}\setminus\mathscr{K} $ such that ${A^L}' = A^L+*$ in normal play. Then we calculate that $o^-({A^L}' +B) = o^+({A^L}' + B + *) =o^+(A^L+*+B+*) = o^+(A^L+B)$, so Left wins going first under mis\`ere play by moving to ${A^L}'$.

The proof that Left wins $G$ going first under normal play if she wins $G$ going first under mis\`ere play is symmetric. If Left wins $G$ going first under mis\`ere play by moving to $G^L \in\ac(  \mathscr{K} \cup \mathscr{B})$, then just as before,  she wins $G$ going first under normal play with the same move. If Left wins $G$ going first under mis\`ere play by moving to $G^L  = A^L + B \in \ac( (\mathscr{A}\setminus \mathscr{K} )\cup \mathscr{B})$, then we calculate that she wins $G$ going first under normal play by moving to ${A^L}' = A^L +*$, since $o^+({A^L}' +B) = o^+(A^L+*+B) = o^-(A^L+*+B+*) = o^-(A^L+B)$. 

    Say that $A \in \mathscr{A} \setminus \mathscr{K}$. Then we need to show that $o^+(G) = o^-(G+*)$ and $o^-(G) = o^+(G+*)$. 

    We first show $o^+(G) = o^-(G+*)$. Say that Left wins $G$ moving first under normal play. If she wins by moving to $G^L \in \ac( (\mathscr{A}\setminus \mathscr{K} )\cup \mathscr{B})$, then she wins $G+*$ under mis\`ere play by moving to $G^L$ as well, since $o^+(G^L) = o^-(G^L+*)$. If Left wins $G$ moving first under normal play by moving to $G^L \in \ac( (\mathscr{A}\setminus \mathscr{K} )\cup \mathscr{B})$, then she either moved to $A^L + B$ with $A^L \in \mathscr{K}$, or $A+B^L$ with $B^L \in \ac( (\mathscr{A}\setminus \mathscr{K} )\cup \mathscr{B})$. But $A^{\mathcal{L}} \cap  \mathscr{K}$ and $B^{\mathcal{L}}  \cap \text{ac}(\mathscr{K} \cup \mathscr{B})$ are both star-closed, so either way there exists ${G^L}'$ such that ${G^L}' = G^L+*$ in normal play. Then $o^+({G^L}') = o^+(G^L+*) = o^-(G^L+*+*) = o^-(G^L)$, where the second equality is justified by the fact that we are inducting on the birthday of the normal play canonical form of $G$. Since $G^L+* = G^{L'}$, which is an option of $G$, the birthday of the normal play canonical form of $G^L +*$ is less than that of $G$, proving the induction is justified.

The other half of the proof is nearly symmetric. Say Left wins $G+*$ going first under mis\`ere play. By assumption, if $G \in \ac((\mathscr{A}\setminus\mathscr{K} )\cup 
\mathscr{B})$, then moving to $G$ is never the unique winning move in $G+*$ under either play convention for either player. So Left may win $G+*$ going first under mis\`ere play by moving to some $G^L+*$. If $G^L \in \ac( (\mathscr{A}\setminus \mathscr{K} )\cup \mathscr{B})$, then $o^-(G^L+*) =o^+(G^L)$, so Left wins $G$ going first in normal play by moving to $G^L$. If $G^L \in \ac(  \mathscr{K} \cup \mathscr{B})$, then there exists ${G^L}' \in\ac(  \mathscr{K} \cup \mathscr{B})$ such that ${G^L}' = G^L+*$. Then $o^+({G^L}') =o^+(G^L+*) = o^-(G^L+*)$, so Left wins $G$ going first under normal play by moving to ${G^L}'$.

The proof that $o^-(G) = o^+(G+*)$ is symmetric.
\end{proof}

\end{document}
